\section{Introduction}

\label{sec:intro}
%% =====================================================================

Financial systems have two fundamental requirements for record-keeping:
\emph{durability} (records must not be lost) and \emph{integrity}
(records must not be silently modified).  Traditional architectures
use a centralized database (PostgreSQL, Oracle, SQL Server) as the
source of truth, protected by access controls, audit logs, and backup
procedures.  The integrity guarantee is institutional: it holds because
the operator is trustworthy and the database is properly administered.

Blockchain provides a stronger guarantee.  Every state transition is
recorded as a digitally signed transaction in an append-only ledger.
The integrity guarantee is cryptographic: modifying any record requires
forging a digital signature or controlling a supermajority of
validators.  The durability guarantee is distributed: the ledger is
replicated across every full node.

The chain-first architecture makes the blockchain the single source of
truth and reduces the database to a materialized view---a read cache
that accelerates queries but contains no original data.  If the cache
is lost, it is fully reconstructable from the chain.  If the chain
and cache disagree, the chain wins.

\subsection{Why Not Database-First?}

In a database-first architecture, the blockchain is a settlement
rail: the database records a trade, and a background job eventually
posts the settlement transaction to the chain.  This creates several
problems:

\begin{enumerate}[nosep]
  \item \textbf{Split brain.}  If the database says Alice owns 100
        shares but the chain says she owns 95, which is correct?  The
        database-first architecture has no deterministic answer.
  \item \textbf{Silent corruption.}  A database administrator can
        modify records without leaving a cryptographic trace.  The
        audit log is stored in the same database it is supposed to
        protect.
  \item \textbf{Backup divergence.}  Database backups taken at
        different times may reflect inconsistent states.  Restoring
        from backup can lose or duplicate settlements.
  \item \textbf{Regulatory risk.}  SEC Rule 17a-4 requires records
        to be stored in non-rewritable, non-erasable format (WORM).
        A mutable database does not satisfy this requirement without
        additional infrastructure.
\end{enumerate}

\subsection{Chain-First: The Inversion}

In the chain-first architecture:

\begin{itemize}[nosep]
  \item The chain is WORM (Write Once, Read Many) by construction.
  \item Every state transition is a signed transaction.
  \item The database is a read cache, always derivable from chain
        events.
  \item If the cache is corrupted or lost, it is rebuilt by replaying
        the chain from genesis.
\end{itemize}

%% =====================================================================
